\subsection{Stenotic Flow Observables and Model Tiers}
\label{subsec:stenotic-hemodynamics-observables}

A stenosis model must declare both its mathematical tier and its observable. A
resolved 3D or FSI model can represent local velocity gradients, wall interaction,
and geometric detail directly, while a distributed 1D model stores axial area and
flow and recovers other quantities through closure and observation maps
\cite{FormaggiaEtAl2007FSICoupling,QuarteroniVeneziani2003GeometricalMultiscale,KopplHelmig2023DimensionReducedModeling}. The present report uses this hierarchy
to motivate a reduced area--flow solver and a velocity comparison operator.

\begin{table}[!htb]
    \centering
    \scriptsize
    \caption{Focused model-tier roles for this report. Each tier requires its own
    model contract and output map before results can be compared.}
    \label{tab:model-tier-review}
    \begin{tabularx}{\textwidth}{@{}
        >{\raggedright\arraybackslash}p{0.20\textwidth}
        >{\raggedright\arraybackslash}p{0.30\textwidth}
        >{\raggedright\arraybackslash}X@{}}
        \toprule
        \textbf{Tier} & \textbf{Retained information} & \textbf{Role in this report} \\
        \midrule
        3D CFD / FSI
            & Resolved velocity, geometry, and possibly wall motion.
            & Source of available resolved velocity data and matching limits, not the
            selected forward solver. \\
        1D distributed
            & Axial area--flow dynamics with wall, profile, friction, rheology, and
            geometry closures.
            & Selected solver tier and main implementation contract. \\
        0D / network and multiscale
            & Lumped or coupled upstream/downstream effects and interface conditions.
            & Boundary and coupling context; not the completed comparison target. \\
        \bottomrule
    \end{tabularx}
\end{table}

\subsection{One-Dimensional Stenosis Models and Closure Dependence}
\label{subsec:model-hierarchy-literature}
\label{subsec:literature-1d-stenotic-vessels}

One-dimensional blood-flow models reduce a vessel to cross-sectional area and
flow while retaining axial propagation and source terms
\cite{FormaggiaEtAl2003OneDimensionalBloodFlow,QuarteroniFormaggia2004CardiovascularModeling,ShiEtAl2011BloodFlowModelReview}. In a stenosis, that reduction is
sensitive to how the narrowed reference geometry enters the balance law.
Classical stenosis-loss formulas and unsteady corrections model pressure-loss
closures, while distributed 1D stenosis models place geometry and wall response
inside an area--flow system
\cite{YoungTsai1973SteadyStenosis,YoungTsai1973UnsteadyStenosis,LyrasLee2021PressureDropROM,CanicEtAl2024Extended1DStenosis}. Network-scale arterial models add
wall and terminal assumptions that must also be declared
\cite{DuanmuEtAl2019CoronaryArterialTree,Wang2014BloodFlowNetworks}.

The literature therefore makes closure dependence central. A momentum-flux factor
changes convective transport; a velocity profile changes both flux and any
radial reconstruction; a wall law changes wave speed and area response; a
friction or rheology law changes dissipation; and a boundary approximation
changes how incoming information is imposed. The selected implementation is an
$R_{\max}$-normalized Canic-derived variant, so the literature role is to demand
an explicit implementation contract rather than to treat all 1D stenosis models
as interchangeable.

\begin{table}[!htb]
    \centering
    \scriptsize
    \caption{Focused literature roles behind the present research gap. The rows
    are not a systematic survey; they identify the evidence categories needed for
    the implemented solver.}
    \label{tab:sota-model-method-map}
    \begin{tabularx}{\textwidth}{@{}
        >{\raggedright\arraybackslash}p{0.23\textwidth}
        >{\raggedright\arraybackslash}p{0.33\textwidth}
        >{\raggedright\arraybackslash}X@{}}
        \toprule
        \textbf{Literature role} & \textbf{Representative support} & \textbf{Implication for this report} \\
        \midrule
        Stenosis-aware 1D models
            & Area--flow reductions with wall, profile, friction, and geometry
            closures
            \cite{FormaggiaEtAl2003OneDimensionalBloodFlow,CanicEtAl2024Extended1DStenosis}.
            & The implemented model must state its coordinate map and closure choices. \\
        Open solver and benchmark practice
            & Reproducible 1D solver workflows and benchmark cases
            \cite{BoileauEtAl2015Benchmark,BenemeritoEtAl2024OpenBF}.
            & Runnable records and persisted diagnostics are part of the evidence. \\
        Well-balanced balance-law methods
            & Source-aware discretizations for vessel equilibria
            \cite{LeVeque2002FiniteVolumeMethodsForHyperbolicProblems,PimentelGarciaEtAl2025HighOrderWellBalanced}.
            & Rest-equilibrium preservation is a primary numerical property, not a
            secondary detail. \\
        Cross-dimensional observables
            & Multiscale and reduced/resolved comparisons requiring compatible
            interface or observation quantities
            \cite{FormaggiaEtAl2007FSICoupling,QuarteroniVeneziani2003GeometricalMultiscale}.
            & A declared plane-quadrature operator is required before comparing 1D and
            resolved 3D velocity outputs. \\
        \bottomrule
    \end{tabularx}
\end{table}

\begin{table}[!htb]
    \centering
    \scriptsize
    \caption{Closure dimensions that should be separated when comparing
    stenosis-aware 1D models.}
    \label{tab:one-dimensional-closure-review}
    \begin{tabularx}{\textwidth}{@{}
        >{\raggedright\arraybackslash}p{0.21\textwidth}
        >{\raggedright\arraybackslash}p{0.31\textwidth}
        >{\raggedright\arraybackslash}X@{}}
        \toprule
        \textbf{Closure dimension} & \textbf{Common choices} & \textbf{Effect on evidence} \\
        \midrule
        State variables
            & Area--flow, radius--flow, pressure--flow, or compartment variables.
            & Determines native outputs and required reconstructions. \\
        Wall model
            & Fixed geometry, algebraic wall law, compliant wall, terminal compliance.
            & Controls area response and wave speed. \\
        Profile / momentum flux
            & Flat, parabolic, power-law, empirical correction factors.
            & Controls convective flux and reconstructed radial profiles. \\
        Friction / rheology
            & Newtonian, Carreau, Carreau--Yasuda, Casson, power-law, tube-flow laws
            \cite{PriesEtAl1992BloodViscosityTubeFlow,KumarEtAl2019CarreauYasudaHemodynamics,LiuEtAl2021ICASRheology}.
            & Controls dissipative closure and shear-rate proxy interpretation. \\
        Stenosis treatment
            & Reference-radius variation, empirical loss terms, variable-radius
            distributed terms.
            & Determines how geometry enters the momentum balance. \\
        Boundary and numerics
            & Characteristic rules, terminal models, finite-volume, DG, WENO,
            Lax--Friedrichs-type, or well-balanced schemes.
            & Determines conservation, equilibrium preservation, stability evidence,
            and boundary-regime diagnostics. \\
        Evidence category
            & Manufactured solutions, equilibrium tests, benchmarks, resolved fields,
            experimental data.
            & Separates implementation verification, diagnostic cross-model comparison,
            and external validation. \\
        \bottomrule
    \end{tabularx}
\end{table}
